\documentclass[11pt,a4paper]{article}

% ============================================================
% Typography and page layout
% ============================================================

\usepackage[T1]{fontenc}
\usepackage[utf8]{inputenc}
\usepackage{lmodern}
\usepackage[a4paper,top=2.2cm,bottom=2.3cm,left=2.35cm,right=2.35cm]{geometry}
\usepackage{microtype}
\usepackage{setspace}
\setstretch{1.10}

\setlength{\parindent}{0pt}
\setlength{\parskip}{0.55em}
\setlength{\emergencystretch}{2em}

% ============================================================
% Mathematics, tables and figures
% ============================================================

\usepackage{amsmath}
\usepackage{amssymb}
\usepackage{booktabs}
\usepackage{tabularx}
\usepackage{array}
\usepackage{longtable}
\usepackage{graphicx}
\usepackage{float}
\usepackage{placeins}
\usepackage{csvsimple}
\usepackage{siunitx}
\usepackage{xstring}


\sisetup{
  round-mode=places,
  round-precision=4,
  detect-all,
  group-digits=false,
  table-number-alignment=center,
  input-symbols = {-,+}
}

\usepackage{enumitem}
\setlist[itemize]{
  leftmargin=1.4em,
  itemsep=0.15em,
  topsep=0.25em
}

% ============================================================
% Captions, colors and hyperlinks
% ============================================================

\usepackage{xcolor}
\definecolor{linkblue}{RGB}{28,72,122}
\definecolor{lightgray}{RGB}{245,246,248}
\definecolor{midgray}{RGB}{90,90,90}

\usepackage[
  font=small,
  labelfont=bf,
  justification=centering,
  singlelinecheck=false,
  skip=6pt
]{caption}

\usepackage[
  colorlinks=true,
  linkcolor=linkblue,
  citecolor=linkblue,
  urlcolor=linkblue,
  pdfauthor={Momen Ismail},
  pdftitle={Machine Learning for Cross-Sectional Stock Return Prediction}
]{hyperref}

\usepackage[nameinlink,noabbrev]{cleveref}

% ============================================================
% Bibliography
% ============================================================

\usepackage{csquotes}
\usepackage[
  backend=biber,
  style=authoryear,
  maxcitenames=2,
  giveninits=true,
  uniquename=init,
  doi=false,
  url=false,
  isbn=false
]{biblatex}

\addbibresource{references.bib}

% ============================================================
% Header and footer
% ============================================================

\usepackage{fancyhdr}
\pagestyle{fancy}
\fancyhf{}
\fancyhead[L]{Machine Learning for Finance}
\fancyhead[R]{Momen Ismail}
\fancyfoot[C]{\thepage}
\renewcommand{\headrulewidth}{0.4pt}
\setlength{\headheight}{14pt}

% ============================================================
% Section formatting
% ============================================================

\usepackage{titlesec}

\titleformat{\section}
  {\large\bfseries}
  {\thesection}
  {0.65em}
  {}

\titleformat{\subsection}
  {\normalsize\bfseries}
  {\thesubsection}
  {0.60em}
  {}

\titlespacing*{\section}{0pt}{1.4em}{0.45em}
\titlespacing*{\subsection}{0pt}{1.0em}{0.25em}

\setcounter{secnumdepth}{2}

% ============================================================
% Paths to existing project outputs
% ============================================================

\newcommand{\ResultsPath}{../output/models/interpretation}
\newcommand{\FigurePath}{../output/models/interpretation/figures}

% ============================================================
% Helper commands
% ============================================================

\ExplSyntaxOn

\NewDocumentCommand{\CSVNumber}{m}
 {
  \tl_if_blank:eTF {#1}
   {--}
   {\exp_args:Ne \num {#1}}
 }

\NewDocumentCommand{\ModelName}{m}
 {
  \exp_args:Ne \str_case:nnF {#1}
   {
    {historical_mean}{Historical~Mean}
    {ols_3}{OLS-3}
    {pls}{PLS}
    {elastic_net}{Elastic~Net}
    {decision_tree}{Decision~Tree}
    {random_forest}{Random~Forest}
   }
   {\texttt{\tl_to_str:n {#1}}}
 }

\NewDocumentCommand{\FamilyName}{m}
 {
  \exp_args:Ne \str_case:nnF {#1}
   {
    {random_forest}{Random~Forest}
    {elastic_net}{Elastic~Net}
    {decision_tree}{Decision~Tree}
    {pls}{PLS}
   }
   {\texttt{\tl_to_str:n {#1}}}
 }

\ExplSyntaxOff

\newcommand{\papertitle}
  {Machine Learning for Cross-Sectional Stock Return Prediction}

% ============================================================
% Document
% ============================================================

\begin{document}

% ============================================================
% Compact title block
% ============================================================

\begin{center}
  {\LARGE\bfseries \papertitle\par}

  \vspace{0.35em}

  {\large\bfseries Preliminary Project Report\par}

  \vspace{0.7em}

  {\large Momen Ismail\par}

  \vspace{0.25em}

  {Machine Learning for Finance, University of Bonn\par}

  {Instructor: Ivan Gufler\par}

  \vspace{0.35em}

  {\today\par}
\end{center}

\vspace{0.6em}
\hrule
\vspace{0.8em}

% ============================================================
% Abstract
% ============================================================

\begin{abstract}
\noindent
This paper compares linear, regularized, dimension-reduction and tree-based
methods for predicting next-month stock excess returns. The models are
evaluated using a chronological design that separates model selection from a
2020--2025 test period. Partial Least Squares produces the lowest preliminary
test-period mean squared error and an out-of-sample \(R^2\) of approximately
\(0.18\%\). Random Forest and a parsimonious three-variable OLS model also
slightly outperform the historical-mean benchmark. Elastic Net selects no
nonzero coefficients and effectively reproduces the benchmark, while the
single Decision Tree performs substantially worse. The results illustrate
that machine-learning flexibility does not automatically translate into
reliable stock-return forecasts and that small improvements must be interpreted
cautiously. These estimates are preliminary: a project audit identified target
continuity, macro-timing and workflow-documentation issues that will be
corrected before a complete clean rebuild.
\end{abstract}

\vspace{0.4em}

\noindent
\textbf{Keywords:}
stock returns, machine learning, Partial Least Squares, Random Forest,
out-of-sample prediction

% ============================================================
% 1. Introduction
% ============================================================

\section{Introduction}
\label{sec:introduction}

Expected stock returns are central to asset pricing, portfolio allocation and
risk management. Predicting them is nevertheless difficult because realized
monthly returns contain substantial idiosyncratic variation and because
predictive relationships may change across market regimes. A model that fits
historical observations well may therefore perform poorly when applied to
future data.

Machine-learning methods are potentially attractive in this setting because
they can process many correlated predictors, impose regularization and capture
nonlinearities or interactions. At the same time, their flexibility increases
the risk of overfitting. A meaningful evaluation must consequently be based on
genuinely out-of-sample observations and must compare complex models with a
simple benchmark.

This paper asks:

\begin{quote}
\emph{Can linear dimension-reduction, regularized regression and tree-based
methods improve next-month stock excess-return forecasts relative to a
historical-mean benchmark?}
\end{quote}

The study compares six approaches under a common information set and evaluation
framework: the historical mean, OLS with three economically motivated
predictors, Partial Least Squares, Elastic Net, a Decision Tree and a Random
Forest. Model selection is performed chronologically, while 2020--2025 is
used as the preliminary test period in the current implementation.

The study is related to the predictive asset-pricing literature, particularly
the use of high-dimensional firm characteristics in
\textcite{gu2020empirical}. It also follows the benchmark-oriented perspective
of \textcite{welch2008comprehensive}, who show that apparently sophisticated
return-prediction models often struggle to outperform simple historical
forecasts.

The contribution is not the introduction of a new algorithm. Instead, the
paper provides a transparent and reproducible comparison of models with
different levels of flexibility. It places particular emphasis on negative
findings, model instability and the distinction between statistical ranking
and economically meaningful predictability.

This document is a preliminary project report. It describes the current
implementation and reports its existing outputs transparently, but it does not
present them as final empirical evidence. The project audit identified a small
number of consequential timing and workflow issues. These are documented below
and will be corrected before the output directory is deleted and the entire
project is rebuilt from the permanent inputs.

% ============================================================
% 2. Data
% ============================================================

\section{Data}
\label{sec:data}

\subsection{Sources and observation unit}

The observation unit is a stock-month. The stock universe is based on a locked
current-and-historical S\&P 500-style ticker list. Daily price and volume
information comes from Yahoo Finance. Market index and VIX series are used to
construct aggregate market and volatility variables. Monthly Fama--French data
provide the risk-free rate used to calculate excess returns. Annual Compustat
data supply accounting characteristics and industry classifications.
Welch--Goyal variables provide aggregate predictive states.

Daily price histories begin before the modeling period so that long-window
characteristics can be constructed. The current monthly modeling sample begins in
January 1990. The held-out test period covers January 2020 through December
2025.

For stock \(i\) at month \(t\), the intended target is the excess return in the
immediately following calendar month:

\begin{equation}
 y_{i,t+1}
 =
 r_{i,t+1}-RF_{t+1},
 \label{eq:target}
\end{equation}

where \(r_{i,t+1}\) is the stock return and \(RF_{t+1}\) is the one-month
risk-free rate.

\subsection{Timing and preprocessing}

Predictors are constructed using information available by the end of month
\(t\). Momentum measures exclude the return being predicted. Annual Compustat
observations are made available six months after fiscal year-end and are merged
backward to the most recent report available for each ticker-month. This
conservative lag reduces, but cannot fully eliminate, the risk of look-ahead
bias because exact filing dates are unavailable.

The project audit confirms that stock-level feature construction and the
six-month Compustat availability rule are largely chronological. It also finds
two timing issues that must be corrected before final estimation. First, the
current ticker-level lead operation generally constructs the intended target,
but in 26 retained observations the next available ticker observation is not
the immediately following calendar month. These nonconsecutive transitions
will be removed. Second, the Welch--Goyal macro variables are currently merged
in the same month. An explicit one-month availability lag will be imposed so
that the final rebuild uses only macro information available at the forecast
date.

Missing predictor observations are replaced using the contemporaneous
cross-sectional median. If an entire predictor is missing within one month, the
remaining values are filled with zero. Predictors are then converted to
within-month cross-sectional ranks on the interval \([-1,1]\):

\begin{equation}
 x^{\mathrm{rank}}_{i,t}
 =
 2
 \frac{\operatorname{rank}(x_{i,t})-1}
      {N_t-1}
 -1.
 \label{eq:rank-transform}
\end{equation}

This transformation reduces sensitivity to units and extreme predictor values
while preserving the relative ordering of firms within each month.

\subsection{Predictor set}

The current design contains 499 predictors. The base variables represent major
economic groups:

\begin{itemize}
  \item size, valuation, profitability, investment and leverage;
  \item momentum, reversals and price trends;
  \item volatility, downside risk, beta and idiosyncratic volatility;
  \item liquidity and trading activity;
  \item industry indicators;
  \item market, VIX and Welch--Goyal aggregate states;
  \item interactions between firm characteristics and aggregate states.
\end{itemize}

The resulting panel is high-dimensional and highly correlated. This provides
the motivation for regularization, dimension reduction and nonlinear methods.

\subsection{Chronological sample design}

The data are divided chronologically:

\begin{itemize}
  \item training period: 1990--2014;
  \item validation period: 2015--2019;
  \item preliminary test period: 2020--2025.
\end{itemize}

Hyperparameters are selected using annual expanding-window folds within the
training period. Fixed and optimized specifications are compared using the
validation period. Once the specifications are selected, the training and
validation observations are combined into a development sample, each model is
refitted once and evaluated on the preliminary test period. The project audit
also identified inconsistencies between the written description of the tuning
periods and selection rule and their implementation. The exact chronology and
selection rule will therefore be reconciled and documented before the clean
rebuild.

Random cross-validation is not used because it would allow future market
observations to influence earlier forecasts. The chronological design more
closely represents the information available to an investor in real time.

% ============================================================
% 3. Methodology
% ============================================================

\section{Methodology}
\label{sec:methodology}

\subsection{Forecasting models}

The general forecasting problem is

\begin{equation}
 y_{i,t+1}
 =
 f(\mathbf{x}_{i,t})
 +
 \varepsilon_{i,t+1},
 \label{eq:forecast}
\end{equation}

where the function \(f(\cdot)\) differs across models.

\textbf{Historical Mean.}
The benchmark predicts the mean target observed in the permitted estimation
sample. It ignores firm characteristics and represents the performance of a
simple constant forecast.

\textbf{OLS-3.}
The parsimonious OLS specification uses size, book-to-market and 12-month
momentum. These variables offer an interpretable reference between the
constant benchmark and the high-dimensional models.

\textbf{Partial Least Squares.}
PLS constructs supervised latent components from the 499 predictors. Unlike
unsupervised dimension reduction, the components are chosen to be related to
the target. PLS is particularly relevant when predictors are numerous and
strongly correlated \parencite{wold1984collinearity}.

\textbf{Elastic Net.}
Elastic Net combines \(L_1\) and \(L_2\) regularization:

\begin{equation}
 \widehat{\boldsymbol{\beta}}
 =
 \arg\min_{\boldsymbol{\beta}}
 \left[
 \frac{1}{2n}
 \left\|
 \mathbf{y}-\mathbf{X}\boldsymbol{\beta}
 \right\|_2^2
 +
 \alpha
 \left(
 \rho\|\boldsymbol{\beta}\|_1
 +
 \frac{1-\rho}{2}
 \|\boldsymbol{\beta}\|_2^2
 \right)
 \right].
 \label{eq:elastic-net}
\end{equation}

The \(L_1\) component permits variable selection, while the \(L_2\) component
stabilizes estimates when predictors are correlated
\parencite{tibshirani1996regression,zou2005regularization}.

\textbf{Decision Tree.}
A regression tree recursively divides the predictor space and assigns a
constant prediction to each terminal leaf. It can represent threshold effects
and interactions but is sensitive to changes in the estimation sample.

\textbf{Random Forest.}
Random Forest averages predictions across bootstrap trees and considers a
random subset of predictors at each split \parencite{breiman2001random}.
Averaging is intended to reduce the variance of an individual tree while
preserving nonlinear predictive relationships.

\subsection{Evaluation metrics}

The primary criterion is monthly mean squared error:

\begin{equation}
 \operatorname{MSE}_{\mathrm{month}}
 =
 \frac{1}{T}
 \sum_{t=1}^{T}
 \left[
 \frac{1}{N_t}
 \sum_{i=1}^{N_t}
 \left(
 y_{i,t+1}
 -
 \widehat{y}_{i,t+1}
 \right)^2
 \right].
 \label{eq:monthly-mse}
\end{equation}

Errors are first averaged across firms and then across months, preventing months
with more observations from receiving greater weight.

The report also presents RMSE, MAE, prediction-target correlation and
out-of-sample \(R^2\):

\begin{equation}
 R^2_{\mathrm{OOS}}
 =
 1
 -
 \frac{
   \sum_{i,t}
   (y_{i,t+1}-\widehat y_{i,t+1})^2
 }{
   \sum_{i,t}
   (y_{i,t+1}-\widehat y^{\mathrm{HM}}_{i,t+1})^2
 }.
 \label{eq:oos-r2}
\end{equation}

A positive value means that the model reduces squared forecast error relative
to the historical mean. A negative value means that the model performs worse
than the benchmark.

% ============================================================
% 4. Results
% ============================================================

\section{Preliminary Results}
\label{sec:results}

\subsection{Overall predictive performance}

\begin{table}[H]
  \centering
  \caption{Preliminary test-period performance, 2020--2025}
  \label{tab:final-results}

  \small
  \setlength{\tabcolsep}{5.5pt}

  \begin{tabular}{
    c
    l
    S[table-format=1.5]
    S[table-format=1.5]
    S[table-format=1.5]
    S[table-format=-1.5]
    S[table-format=-1.5]
  }
    \toprule
    {Rank}
    & {Model}
    & {MSE}
    & {RMSE}
    & {MAE}
    & {\(R^2_{\mathrm{OOS}}\)}
    & {Correlation} \\
    \midrule

    \csvreader[
      head to column names,
      late after line=\\
    ]{\ResultsPath/final_prediction_results.csv}{
      rank=\Rank,
      model=\Model,
      monthly_mse=\MonthlyMSE,
      rmse=\RMSE,
      mae=\MAE,
      oos_r2=\OOSR,
      correlation=\Correlation
    }{
      \Rank
      & \ModelName{\Model}
      & \CSVNumber{\MonthlyMSE}
      & \CSVNumber{\RMSE}
      & \CSVNumber{\MAE}
      & \CSVNumber{\OOSR}
      & \CSVNumber{\Correlation}
    }

    \bottomrule
  \end{tabular}

  \vspace{0.35em}

  \begin{minipage}{0.92\textwidth}
    \footnotesize
    \textit{Notes:}
    Models are ranked by monthly MSE. Correlation is undefined for constant or
    near-constant predictions. Positive \(R^2_{\mathrm{OOS}}\) indicates an
    improvement over the historical-mean benchmark.
  \end{minipage}
\end{table}

PLS provides the lowest monthly MSE and the highest out-of-sample \(R^2\).
Its \(R^2_{\mathrm{OOS}}\) of \(0.00179\) means that it reduces squared forecast
error by approximately \(0.18\%\) relative to the historical mean. The result
is positive, but economically small.

Random Forest ranks second with an out-of-sample \(R^2\) of approximately
\(0.058\%\). OLS-3 ranks third with approximately \(0.049\%\). The difference
between these two models is minimal, indicating that the large increase in
model complexity delivered by Random Forest produces only a small additional
reduction in squared error relative to the parsimonious OLS specification.

Elastic Net and the historical mean produce almost identical results. The
optimized Elastic Net contains no coefficients above the specified tolerance
and therefore behaves as a constant predictor. This is a meaningful negative
finding: after chronological validation, the regularized model found no
sufficiently stable incremental signal among the 499 predictors.

The Decision Tree performs considerably worse than all other models. Its
out-of-sample \(R^2\) is approximately \(-13.3\%\), and its prediction-target
correlation is negative. The individual tree appears to fit unstable splits
that do not generalize to the test sample.

\begin{figure}[H]
  \centering
  \includegraphics[
    width=0.86\textwidth,
    trim=5 5 5 22,
    clip
  ]{\FigurePath/test_oos_r2_by_model.png}

  \caption{Preliminary test-period out-of-sample \(R^2\) by model}
  \label{fig:test-oos-r2}
\end{figure}

\Cref{fig:test-oos-r2} emphasizes both the ranking and the limited absolute
predictive power. PLS, Random Forest and OLS-3 outperform the benchmark, but all
positive values remain close to zero. The Decision Tree is qualitatively
different because its additional flexibility substantially increases forecast
error.

The MAE ranking differs slightly from the MSE ranking. The historical mean and
Elastic Net have the lowest MAE, while PLS has the lowest MSE. This suggests
that PLS reduces some relatively large forecast errors, which receive more
weight under squared loss, without reducing the typical absolute error.

\subsection{Financial interpretation}

The small positive out-of-sample \(R^2\) values represent small reductions in
forecast error relative to the historical-mean benchmark. They provide evidence
of weak statistical predictability in the current project version, but they do
not establish profitable trading. In particular, MSE does not account for
portfolio turnover, transaction costs, liquidity, short-selling constraints or
risk adjustment. Nor does it show that the cross-sectional ordering of forecasts
can be converted into stable implementable positions. The present evidence
should therefore be interpreted as a preliminary comparison of forecasting
methods, not as a ready-to-use investment strategy.

\subsection{Performance across test years}

The preliminary pooled ranking does not imply stable performance in every year. PLS and
Random Forest achieve positive yearly out-of-sample \(R^2\) in half of the six
test years. OLS-3 does so in one third of the years. The Decision Tree beats the
benchmark in only one year.

\begin{figure}[H]
  \centering
  \includegraphics[
    width=0.90\textwidth,
    trim=5 5 5 22,
    clip
  ]{\FigurePath/yearly_oos_r2_by_model.png}

  \caption{Out-of-sample \(R^2\) by model and test year}
  \label{fig:yearly-oos-r2}
\end{figure}

The yearly results demonstrate that even the best overall model does not
consistently outperform the historical mean. This instability is important
because the test period contains only six years and includes unusually
different market environments. The overall advantage of PLS should therefore
be interpreted as an average test-period result rather than evidence of stable
performance in every regime.

\subsection{Forecast dispersion}

\begin{figure}[H]
  \centering
  \includegraphics[
    width=0.90\textwidth,
    trim=5 5 5 22,
    clip
  ]{\FigurePath/prediction_distribution_by_model.png}

  \caption{Distribution of preliminary test-period predictions}
  \label{fig:prediction-distribution}
\end{figure}

The prediction distributions help explain the model results.
Elastic Net is absent as a meaningful distribution because it produces
constant or near-constant forecasts. PLS and OLS generate relatively
concentrated predictions, while the tree-based models produce more dispersed
forecasts. Dispersion alone does not indicate predictive quality. The Decision
Tree produces variable forecasts but performs poorly because those forecasts
are systematically misaligned with realized returns.

\subsection{Feature importance}

\begin{figure}[H]
  \centering
  \includegraphics[
    width=0.88\textwidth,
    trim=5 5 5 5,
    clip
  ]{\FigurePath/top_pls_coefficients.png}

  \caption{Largest PLS coefficients by absolute magnitude}
  \label{fig:pls-coefficients}
\end{figure}

PLS coefficients measure conditional predictive associations after
preprocessing and scaling. Their magnitude should not be interpreted as a
causal effect or as an estimated risk premium. Because PLS combines information
through latent components, the coefficients reflect the joint high-dimensional
prediction problem rather than isolated univariate relationships.

\begin{figure}[H]
  \centering
  \includegraphics[
    width=0.88\textwidth,
    trim=5 5 5 5,
    clip
  ]{\FigurePath/top_random_forest_features.png}

  \caption{Top Random Forest predictors based on impurity importance}
  \label{fig:rf-importance}
\end{figure}

Random Forest importance records how strongly variables reduce squared error
inside the fitted trees. It is useful for describing where the model obtains
its predictive variation, but it has important limitations. Correlated
variables can divide importance among themselves, and impurity importance can
favor predictors with more available split points. The results should therefore
be interpreted at the level of broad predictor groups rather than as a precise
ranking of causal variables.

% ============================================================
% 5. Critical reflection
% ============================================================

\section{Critical Reflection}
\label{sec:critical-reflection}

\subsection{Predictive gains are small}

The main limitation is the weak absolute predictive performance. Although PLS
ranks first, its out-of-sample \(R^2\) is only approximately \(0.18\%\). Random
Forest and OLS-3 improve on the benchmark by even less. Most variation in
monthly stock returns therefore remains unexplained.

A ranking can still be statistically informative, but it should not be confused
with strong forecasting ability. The results do not establish that the small
differences between the leading models would remain stable in another period.

\subsection{Flexible models do not automatically perform better}

The results provide two important failures. First, the individual Decision Tree
performs much worse than the historical mean. Its threshold splits are
apparently too unstable for noisy monthly return data. Random Forest improves
substantially by averaging across trees, illustrating the benefit of variance
reduction.

Second, Elastic Net selects no nonzero coefficients. This does not imply a
software error. Given the selected tuning parameters and chronological
validation procedure, the model preferred complete shrinkage over unstable
predictor effects. The result suggests that individual linear signals are weak
relative to the noise in the target.

\subsection{Implementation issues identified during the project audit}

The project audit identified six implementation and reproducibility problems
that prevent the current outputs from being described as final results:

\begin{itemize}
  \item The ticker-level target lead contains 26 nonconsecutive transitions in
        which the next available observation is not the immediately following
        calendar month.
  \item Welch--Goyal macro variables are merged in the same month rather than
        with the intended explicit one-month availability lag.
  \item The documentation and implementation are inconsistent about the tuning
        periods and the model-selection rule.
  \item The top-level \texttt{run\_pipeline.sh} became stale after the model
        code was reorganized and references stages that no longer match the
        active workflow.
  \item The output structure contains excessive duplicate and legacy artifacts,
        which makes it harder to identify the authoritative results.
  \item Software dependencies are not pinned to exact versions, limiting exact
        reproducibility across environments.
\end{itemize}

These issues do not make the current tables useless: they remain informative
about model behavior and expose several robust negative findings. They do mean,
however, that the numerical rankings must be treated as preliminary. Before the
complete rebuild, the target continuity rule and macro lag will be corrected,
the tuning design and documentation will be reconciled, the pipeline entry
point will be repaired, dependencies will be pinned, and authoritative outputs
will be separated from legacy artifacts. The existing output folder will then
be deleted and regenerated from the permanent inputs in one clean chronological
run.

\subsection{Data limitations}

The stock universe is not a fully point-in-time CRSP constituent universe.
Ticker-based matching may be affected by symbol changes, mergers and incomplete
historical membership records. Survivorship and universe-selection bias cannot
be ruled out completely.

Yahoo Finance does not provide the security master, delisting returns and
microstructure information available in CRSP. Adjusted price histories may also
be revised. Compustat fundamentals are lagged by six months, but exact public
filing dates are not used. The resulting timing design is conservative, though
not perfectly point in time.

The use of annual Compustat information also means that accounting variables
change less frequently than market variables. Some predictors are proxies for
variables used in the published asset-pricing literature and may not reproduce
their original definitions exactly.

\subsection{Validation and specification risk}

Annual expanding-window validation respects chronology and reduces direct
leakage. Nevertheless, the project contains many researcher choices:
predictor construction, interaction design, preprocessing, tuning grids and
sample dates. These decisions may adapt indirectly to known characteristics of
the data.

An untouched holdout period remains credible only if no model changes are
made after inspecting its performance. The current 2020--2025 outputs have
already been examined and are therefore reported only as preliminary evidence.
The corrected workflow must be specified before the clean rebuild, and any
subsequent specification search would require a genuinely new holdout period.

\subsection{No economic trading evaluation}

The project evaluates prediction error but does not construct portfolios.
Consequently, it does not measure turnover, transaction costs, market impact,
short-sale constraints or risk-adjusted investment returns. A small reduction
in MSE cannot automatically be interpreted as an economically profitable
strategy.

A useful extension would specify a portfolio rule before opening a new test
sample. The evaluation could then report gross and net returns, volatility,
Sharpe ratios, turnover and maximum drawdown. That analysis would have to be
kept separate from the current test period to avoid using the same observations
for model assessment and strategy design.

\subsection{What I would do differently}

With improved data access and additional time, the study could be strengthened
in four ways:

\begin{itemize}
  \item use a point-in-time CRSP/Compustat universe with permanent identifiers
        and delisting returns;
  \item use exact filing dates for accounting information;
  \item retain a longer untouched test period or evaluate several historical
        forecast origins;
  \item conduct a pre-specified portfolio evaluation with realistic trading
        costs.
\end{itemize}

Interpretability could also be improved by examining feature groups, repeated
time-block importance and stability across different expanding windows rather
than relying on a single definitive importance ranking.

% ============================================================
% 6. Conclusion
% ============================================================

\section{Conclusion}
\label{sec:conclusion}

This preliminary report examines whether machine-learning methods improve next-month stock
excess-return prediction relative to a historical-mean benchmark. The models
are selected chronologically and evaluated on a held-out 2020--2025 test
period in the current implementation.

PLS performs best according to monthly mean squared error and produces an
out-of-sample \(R^2\) of approximately \(0.18\%\). Random Forest and OLS-3 also
slightly outperform the benchmark. Elastic Net selects no predictors and
effectively collapses to the historical mean. The Decision Tree performs
substantially worse, showing that flexibility without sufficient variance
control can be harmful.

The research question can therefore be answered cautiously. Some statistical
and machine-learning methods improve on the historical mean, but the gains are
small and unstable across years. The results do not demonstrate strong return
predictability or a profitable investment strategy. Their principal lesson is
that, in noisy financial prediction problems, disciplined chronology, simple
benchmarks and honest reporting of failures are at least as important as model
complexity. Because the audit identified target-continuity, macro-timing and
workflow inconsistencies, these results are not the final project evidence.
Those issues will be corrected before the output folder is deleted and the
project is rebuilt cleanly from its permanent inputs.

% ============================================================
% References
% ============================================================

\FloatBarrier
\printbibliography[heading=bibintoc,title={References}]

% ============================================================
% Appendix
% ============================================================

\clearpage
\appendix

\section{Code and Reproducibility Overview}
\label{app:code-overview}

The repository implements the empirical workflow as a sequence of auditable
stages rather than embedding source code in this report:

\begin{enumerate}[leftmargin=1.6em,itemsep=0.2em,topsep=0.25em]
  \item acquisition of permanent external inputs and cleaning of daily stock
        histories;
  \item construction of monthly stock features and the next-month target;
  \item merging of lagged fundamentals and aggregate macroeconomic variables;
  \item final missing-value treatment and within-month rank normalization;
  \item chronological hyperparameter tuning with expanding estimation windows;
  \item preliminary test-period refitting and evaluation;
  \item interpretation, figure generation and report production.
\end{enumerate}

The complete Python and shell code is provided in the accompanying repository.
This document reads its principal tables and figures from the existing
interpretation outputs and does not retrain models or regenerate predictions
during compilation. Before the final clean rebuild, the audited timing,
workflow and dependency issues described in
\cref{sec:critical-reflection} will be corrected.

\section{Additional Results}
\label{app:additional-results}

\subsection{Model complexity}

\begin{table}[H]
  \centering
  \caption{Current model complexity}
  \label{tab:model-complexity}

  \small
  \setlength{\tabcolsep}{5pt}

  \begin{tabular}{
    l
    S[table-format=3.0]
    S[table-format=3.0]
    S[table-format=2.0]
    S[table-format=2.0]
    S[table-format=3.0]
  }
    \toprule
    {Model}
    & {Predictors}
    & {Nonzero}
    & {PLS comp.}
    & {Tree depth}
    & {Forest trees} \\
    \midrule

    \csvreader[
      head to column names,
      late after line=\\
    ]{\ResultsPath/final_model_complexity.csv}{
      model=\Model,
      number_of_predictors=\Predictors,
      nonzero_coefficients=\Nonzero,
      pls_components=\Components,
      tree_depth=\TreeDepth,
      forest_trees=\ForestTrees
    }{
      \ModelName{\Model}
      & \CSVNumber{\Predictors}
      & \CSVNumber{\Nonzero}
      & \CSVNumber{\Components}
      & \CSVNumber{\TreeDepth}
      & \CSVNumber{\ForestTrees}
    }

    \bottomrule
  \end{tabular}
\end{table}

\subsection{Fixed versus optimized specifications}

\begin{table}[H]
  \centering
  \caption{Validation performance of fixed and optimized specifications}
  \label{tab:fixed-optimized}

  \small

  \begin{tabular}{
    l
    S[table-format=1.5]
    S[table-format=1.5]
    S[table-format=-1.5]
    S[table-format=-1.5]
  }
    \toprule
    {Model family}
    & {Fixed MSE}
    & {Optimized MSE}
    & {MSE improvement}
    & {\(\Delta R^2_{\mathrm{OOS}}\)} \\
    \midrule

    \csvreader[
      head to column names,
      late after line=\\
    ]{\ResultsPath/fixed_vs_optimized_results.csv}{
      model_family=\Family,
      fixed_monthly_mse=\FixedMSE,
      optimized_monthly_mse=\OptimizedMSE,
      monthly_mse_improvement=\Improvement,
      oos_r2_change=\RChange
    }{

    \FamilyName{\Family}
    & \CSVNumber{\FixedMSE}
    & \CSVNumber{\OptimizedMSE}
    & \CSVNumber{\Improvement}
    & \CSVNumber{\RChange}
  }

    \bottomrule
  \end{tabular}
\end{table}

\subsection{Additional tree interpretation}

\begin{figure}[H]
  \centering
  \includegraphics[
    width=0.88\textwidth,
    trim=5 5 5 5,
    clip
  ]{\FigurePath/top_decision_tree_features.png}

  \caption{Top Decision Tree predictors based on impurity importance}
  \label{fig:dt-importance}
\end{figure}

The Decision Tree importance figure should be interpreted cautiously because
the model has strongly negative out-of-sample performance. A variable can be
important to a poorly performing fitted model without containing stable
predictive information.

\end{document}
